\documentclass[10pt,twocolumn]{article}
\usepackage[T1]{fontenc}
\usepackage[utf8]{inputenc}
\usepackage{lmodern}
\usepackage[margin=1.8cm,top=2.4cm,bottom=2cm,columnsep=0.6cm]{geometry}
\usepackage{fancyhdr}
\usepackage{titlesec}
\usepackage{booktabs}
\usepackage{amsmath,amssymb,amsfonts}
\usepackage{microtype}
\usepackage{xcolor}
\usepackage{mdframed}
\usepackage{parskip}
\usepackage{enumitem}
\usepackage{hyperref}

\definecolor{rulered}{RGB}{180,30,30}
\definecolor{boxgray}{RGB}{245,245,245}
\definecolor{keywordgray}{RGB}{100,100,100}

\pagestyle{fancy}
\fancyhf{}
\fancyhead[L]{\textsc{\small Claude Mag}}
\fancyhead[C]{\small\textit{Machine Learning Research Digest}}
\fancyhead[R]{\small 2026-09-02}
\fancyfoot[C]{\small\thepage}
\renewcommand{\headrulewidth}{0.8pt}

\titleformat{\section}{\normalsize\bfseries\color{rulered}}{}{0pt}{}%
  [\vspace{-4pt}\color{rulered}\rule{\columnwidth}{0.4pt}]
\titlespacing{\section}{0pt}{8pt}{4pt}

\hypersetup{colorlinks=true,urlcolor=rulered,linkcolor=black}

\begin{document}
\setlength{\parindent}{0pt}
\setlength{\parskip}{4pt}

{\small\color{keywordgray}\textbf{Keywords:}
  constrained decoding \textbullet{}
  context-free grammar \textbullet{}
  structured generation \textbullet{}
  pushdown automata}\par\vspace{4pt}

{\LARGE\bfseries\raggedright
  Stay Within Your Bounds: Distance-Guided
  Decoding for Guaranteed Context-Free
  Grammar Compliance}\par\vspace{4pt}

{\large\itshape\raggedright
  Lookahead Distance to Acceptance Closes the
  Infeasibility Gap in Grammar-Constrained LM Decoding}\par\vspace{6pt}

{\small\textbf{By} Vincenzo Collura, Karim Tit,
  Eleonora Giunchiglia, Mike Papadakis, Maxime Cordy}\par\vspace{2pt}

{\small\color{keywordgray}arXiv:2608.28229 \textbullet{}
  EMNLP 2026 Findings (Long Paper)}
\par\vspace{2pt}\noindent{\color{rulered}\rule{\columnwidth}{0.6pt}}\vspace{6pt}

Grammar-constrained decoding guarantees that a language model's output
satisfies a target context-free grammar—enabling reliable generation of
JSON, SQL, and logical formulas. But existing methods fail silently:
they allow prefixes that no valid completion can extend within the token
budget. This paper fixes that gap with a lookahead-guided framework
that computes upper-bound distances to the nearest accepting PDA state,
uses them to prune infeasible tokens online, and guides beam search
toward acceptance. Every output is guaranteed syntactically valid;
completion quality improves across all evaluated structure types.

\begin{mdframed}[backgroundcolor=boxgray,linecolor=rulered,linewidth=1pt,
    innerleftmargin=6pt,innerrightmargin=6pt,
    innertopmargin=6pt,innerbottommargin=6pt]
\textbf{\small AT A GLANCE}\par\vspace{4pt}
\begin{description}[leftmargin=0pt,labelsep=4pt,itemsep=2pt,parsep=0pt,
    font=\small\bfseries]
  \item[Problem:] Grammar masks block individually illegal tokens but
    permit prefixes from which no budget-feasible completion satisfies
    the grammar — producing invalid outputs despite the mask.
  \item[Why it matters:] Hard syntactic guarantees are essential for
    code generation, structured data extraction, and symbolic planning.
  \item[Method:] Precompute bounded PDA summaries with upper-bound
    distances to acceptance; prune tokens that exceed the remaining
    budget; bias beam search toward acceptance.
  \item[Result:] 100\% syntactic validity on JSON, SQL, and LTL
    benchmarks; improved completion quality over prior constrained
    baselines.
  \item[Evidence:] Experiments across three structure types and
    multiple grammar sizes; ablations on budget sensitivity and beam width.
\end{description}
\end{mdframed}

\section{The Big Picture}

Structured generation tasks require outputs that are syntactically valid
instances of a target language: JSON objects, SQL queries, Python
expressions, or formulas in linear temporal logic (LTL). Language models
produce token sequences autoregressively, and without constraints they
generate syntactically invalid outputs at a significant rate.

Grammar-constrained decoding solves this by tracking a pushdown automaton
(PDA) derived from the target context-free grammar (CFG). At each step, it
masks out tokens that are not consistent with the current PDA state,
ensuring only locally valid tokens are sampled. This prevents any single
token from violating the grammar.

But local validity is not global validity. A sequence of locally valid
tokens can lead the PDA into a configuration from which no complete valid
string can be generated within the remaining budget. The decoder emits a
prefix that is syntactically valid so far, but cannot be completed.

\section{Why Existing Approaches Fall Short}

Two failure modes arise in prior grammar-constrained decoding:

\textbf{Tokenizer-grammar mismatch:} LM subword tokenizers may split
strings at positions that do not align with grammar terminal boundaries.
A token that spans multiple grammar symbols can leave the PDA in a
non-standard mid-rule state that requires specific continuations — which
the remaining token vocabulary may not support.

\textbf{Budget exhaustion:} Even with perfect tokenizer alignment, the
PDA may enter a configuration requiring many more tokens to reach
acceptance than remain in the budget. The decoder continues generating
but cannot produce a valid closing sequence.

Both failure modes share the same root: the token mask checks local
legality (is this token consistent with the current PDA state?) but not
global feasibility (is there any valid completion within budget?).

\section{Core Mechanism}

The paper introduces a \textbf{distance oracle}: for each PDA configuration
$c = (q, \gamma)$ (state $q$, stack content $\gamma$), precompute an
upper-bound estimate $\hat{d}(c)$ of the minimum number of additional
tokens needed to reach any accepting configuration.

At each decoding step with $B$ tokens remaining, a token $a$ is permitted
only if:
\begin{enumerate}[itemsep=2pt,parsep=0pt]
  \item It is locally consistent with the current PDA state (existing mask)
  \item The successor configuration $c' = \delta(c, a)$ satisfies
    $\hat{d}(c') \leq B$
\end{enumerate}

Condition 2 is the new feasibility filter. It ensures that after token $a$
is emitted, the remaining budget is sufficient to reach acceptance.

Among the feasible tokens, beam search is guided by a distance-adjusted
score:
\[
  s'(a) = \log p_\theta(a \mid x_{<t}) - \lambda \cdot \hat{d}(\delta(c, a))
\]
which biases the beam toward tokens that reduce the distance to acceptance,
improving completion quality.

\section{Technical Formulation}

Let $\mathcal{A} = (Q, \Sigma, \Gamma, \delta, q_0, F)$ be the PDA for
the target CFG. Define the \textbf{distance to acceptance}:
\[
  d(c) = \min\!\left\{n \in \mathbb{N} : \exists w \in \Sigma^n,\,
    c \xrightarrow{w} (q_f, \varepsilon),\, q_f \in F\right\}
\]
with $d(c) = \infty$ if no such $w$ exists.

The paper computes upper bounds $\hat{d}(c) \geq d(c)$ by summarizing
PDA reachability with bounded lookahead:
\begin{enumerate}[itemsep=2pt,parsep=0pt]
  \item For each PDA state $q$, compute the shortest path in the
    non-deterministic reachability graph over $(q, \gamma)$ pairs to
    any $q_f \in F$, with stack operations truncated at a bounded depth.
  \item Propagate distances bottom-up through grammar rule bodies.
\end{enumerate}
The computation is done offline before decoding. Online, only a table
lookup per (configuration, token) pair is required.

For a token budget $T$ and a model generating tokens $a_1, \ldots, a_T$:
\[
  a_t^* = \arg\max_{a:\, \hat{d}(\delta(c_t, a)) \leq T - t}
           s'(a)
\]
guarantees that every emitted prefix can be completed within the remaining
budget.

\section{Learning or Inference Procedure}

The framework operates purely at inference time:
\begin{enumerate}[itemsep=2pt,parsep=0pt]
  \item Given a target CFG, compile it to a PDA and compute the distance
    oracle $\hat{d}(\cdot)$ offline.
  \item At each decoding step, compute the set of feasible tokens
    (locally valid and budget-feasible).
  \item Apply beam search with distance-guided scores.
  \item Return the highest-scoring complete parse at the end.
\end{enumerate}

No fine-tuning or modification of the base language model is required.
The framework integrates with any autoregressive LM.

\section{What the Guarantee Says}

The decoder is \textbf{sound}: every generated sequence is accepted by
the target grammar. Specifically, if the initial budget $T$ satisfies
$T \geq \hat{d}(q_0, \bot)$ (the oracle distance from the initial
configuration), then at least one feasible token exists at every step
and the decoder terminates with a valid output.

The distance oracle guarantee is: for every configuration $c$ encountered
during decoding, $\hat{d}(c) \geq d(c)$, so the feasibility filter never
incorrectly prunes a token that leads to a valid completion.

\section{Experimental Findings}

Evaluated on three structure types with three LLMs (GPT-2-XL,
LLaMA-3-8B, Mistral-7B-v0.3):

\begin{itemize}[itemsep=2pt,parsep=0pt]
  \item \textbf{JSON:} 100\% syntactic validity for all models and
    grammar sizes; schema match rate $+9.3$ pp over greedy and $+6.1$ pp
    over token-mask-only baselines.
  \item \textbf{SQL:} 100\% parse validity; execution accuracy $+7.8$ pp
    over token-mask-only decoding on Spider benchmark.
  \item \textbf{LTL:} 100\% formula well-formedness; satisfaction rate
    of generated plans $+11.2$ pp over prior constrained baselines.
  \item \textbf{Latency:} Oracle precomputation takes 0.3--1.8 s for
    practical grammar sizes; online overhead per step is negligible.
\end{itemize}

\section{Ablations and Interpretation}

\textbf{Budget sensitivity:} With tight budgets ($T$ only 10\% above
$d(q_0, \bot)$), the feasibility filter prunes aggressively and
occasionally leaves no feasible token at mid-sequence steps. A
budget cushion of 30\% eliminates this; the paper recommends budget
estimation from grammar depth.

\textbf{Distance-guidance weight $\lambda$:} $\lambda = 0$ (pure
language model score) recovers the token-mask baseline; $\lambda = 0.5$
achieves the best completion quality. Larger $\lambda$ over-constrains
generation toward acceptance and reduces fluency.

\textbf{Beam width:} Distance-guided beam search with width 4
outperforms greedy distance-guided decoding by $+2.1$ pp on JSON
schema match, with diminishing returns beyond width 8.

\textbf{Grammar size:} The oracle scales well to grammars with up to
500 production rules, covering real SQL dialects. For larger grammars,
the bounded-lookahead approximation introduces small oracle errors but
validity is preserved.

\section*{Reference}

Vincenzo Collura, Karim Tit, Eleonora Giunchiglia, Mike Papadakis,
Maxime Cordy.
\textbf{Stay Within Your Bounds: Distance-Guided Decoding for
Guaranteed Context-Free Grammar Compliance.}
arXiv:2608.28229, EMNLP 2026 Findings.
\url{https://arxiv.org/abs/2608.28229}

\end{document}
